\section{Related Work}

\subsection{Paper-derived Reproducibility}

% 
% 论文衍生复现研究关注智能体能否将论文转化为可执行、可测试的研究对象。面向完整复现的基准包括 PaperBench、AutoReproduce、RExBench 和 REPRO-Bench，分别关注代码重建、论文谱系辅助的实验复现、论文方法扩展实现以及可复现性判断 \cite{starace2025paperbench,zhao2026autoreproduce,edwards2026rexbench,hu2025reprobench}。另一类工作关注论文知识的表示与制品构建，包括将代码片段和技术洞见组织为可执行知识图、将论文转换为基于 MCP 的交互式智能体、构建可复用工具集合，以及从异构科学资源中编译智能体技能 \cite{luo2026xkg,miao2025paper2agent,ma2025reusabletoolsets,shen2026skillfoundry,pan2026anything2skill}。这些研究的输出主要是完整实现、可执行制品、执行轨迹或总体复现分数，研究目标是构建或评估完整复现对象。SkillAudit 的验证对象不同：它不负责构建完整复现制品，而是判断一个固定的简化说明能否在冻结的任务边界下替代一个已经构建的完整说明。
%
% 
Paper-derived reproducibility asks whether an agent can turn a research paper into an executable and testable research object. Reconstruction benchmarks such as PaperBench, AutoReproduce, RExBench, and REPRO-Bench study code reconstruction, paper-lineage-assisted reproduction, research-extension implementation, and reproducibility assessment, respectively \cite{starace2025paperbench,zhao2026autoreproduce,edwards2026rexbench,hu2025reprobench}. Representation and construction systems instead organize paper knowledge into executable knowledge graphs, MCP-based interactive agents, reusable toolsets, or skills compiled from heterogeneous scientific resources \cite{luo2026xkg,miao2025paper2agent,ma2025reusabletoolsets,shen2026skillfoundry,pan2026anything2skill}. Their primary outputs are complete implementations, executable artifacts, execution traces, or aggregate reproduction scores, and their goal is to construct or evaluate complete reproduction objects. SkillAudit has a different validation target: it does not construct a complete reproduction artifact, but tests whether a fixed simplified specification can substitute for an already-constructed full specification under a frozen task boundary.

\subsection{Context Compression}

%
% 上下文压缩研究关注如何缩短模型实际接收的文本，同时保留下游行为。已有工作覆盖基于 token、问题感知和句子级的压缩方法，包括 LLMLingua、LongLLMLingua、LLMLingua-2、上下文感知提示压缩、PromptOptMe 和 Perception Compressor \cite{jiang2023llmlingua,jiang2024longllmlingua,pan2024llmlingua2,liskavets2025contextcompression,larionov2025promptoptme,tang2025perception}。相关的技能和执行效率研究进一步将这一问题扩展到可复用智能体说明，研究技能 token 优化、压缩因素分解和工具选择成本 \cite{gao2026skillreducer,xu2026compression,jia2026autotool}。这些方法通常通过长度、延迟、成本或总体任务性能比较压缩策略或技能制品集合，但不会同时记录删除的原文片段、验证程序步骤的依赖完整性，并为某一个固定候选给出具体的接受或拒绝依据。SkillAudit 与这些方法的关系是互补的：压缩方法负责生成候选，而 SkillAudit 负责验证候选是否仍然足以完成指定任务。
%
% 
Context compression asks how model-facing text can be shortened while preserving downstream behavior. Existing methods cover token-level, question-aware, and sentence-level compression, including LLMLingua, LongLLMLingua, LLMLingua-2, context-aware prompt compression, PromptOptMe, and Perception Compressor \cite{jiang2023llmlingua,jiang2024longllmlingua,pan2024llmlingua2,liskavets2025contextcompression,larionov2025promptoptme,tang2025perception}. Related work on reusable agent instructions and execution efficiency extends the same concern to skill-token optimization, controlled decomposition of compression effects, and tool-selection cost \cite{gao2026skillreducer,xu2026compression,jia2026autotool}. These methods mainly compare compression strategies or skill populations using aggregate length, latency, cost, or task-performance measurements, but they do not jointly expose deleted source passages, verify procedural dependency integrity, and justify the decision for one fixed candidate. SkillAudit is complementary to these methods: upstream compression systems produce candidates, whereas SkillAudit determines whether a candidate remains sufficient for a specified task.

\subsection{Agent Evaluation}


% 智能体评测研究关注完整工具使用智能体在不同接口、任务和难度条件下的行为。工具学习与工具构建工作包括 Toolformer、CITI、KTCE 和 AutoTool，分别涉及工具调用能力、工具学习中的通用能力保持、可复用工具集合以及工具选择效率 \cite{schick2023toolformer,hao2025citi,ma2025reusabletoolsets,jia2026autotool}。评测和诊断研究则包括 AgentBoard、MCP-AgentBench、Beyond Accuracy 以及大规模智能体技能评测框架，重点分析中间执行进展、协议化工具交互、能力边界和完整技能效用 \cite{ma2024agentboard,guo2026mcpagentbench,wang2026beyondaccuracy,shaposhnikov2026agenticskills}。这些研究的输出通常是完整智能体的分数、执行轨迹、能力画像或成本估计，而不是某个简化技能说明能否替代完整说明的判断。SkillAudit 将评测单位从完整智能体改变为候选说明，并将三条件执行结果转化为带有来源、运行有效性和判断依据的 Accept、Reject 或 Invalid 结论。
%
% 
Agent evaluation asks how complete tool-using agents behave across interfaces, tasks, and difficulty regimes. Tool-learning and tool-construction studies include Toolformer, CITI, KTCE, and AutoTool, covering tool invocation, preservation of general ability during tool learning, reusable toolsets, and tool-selection efficiency \cite{schick2023toolformer,hao2025citi,ma2025reusabletoolsets,jia2026autotool}. Evaluation and diagnostic studies, including AgentBoard, MCP-AgentBench, Beyond Accuracy, and large-scale agent-skill evaluation frameworks, analyze intermediate progress, protocol-aware tool interaction, capability boundaries, and whole-skill utility \cite{ma2024agentboard,guo2026mcpagentbench,wang2026beyondaccuracy,shaposhnikov2026agenticskills}. Their outputs are generally scores, execution traces, capability profiles, or cost estimates for complete agents, rather than decisions about whether a simplified skill specification can replace a full one. SkillAudit complements these evaluations by changing the unit of analysis from a complete agent to a candidate specification and by converting three-condition execution results into a source-linked \textsc{Accept}, \textsc{Reject}, or \textsc{Invalid} decision.